\section{Finite Resource Sampling: Torus Translation of Internal Phase and Periodic Approximation}

\subsection{Internal Phase Dynamics and One-Dimensional Factor}

\begin{definition}[Internal phase translation]\label{def:internal-phase-translation}
Let the quotient space of internal space be
\[
\TT_{\perp}:=E_{\perp}/\pi_{\perp}(L)
\]
a compact abelian group (typically a three-dimensional torus). Take a translation vector $\beta\in \TT_{\perp}$. Define the discrete-time dynamics
\[
u_{t+1}=u_t+\beta,\qquad t\in\ZZ_{\ge 0}.
\]
\end{definition}

\begin{definition}[One-dimensional observation factor]\label{def:one-dimensional-factor}
Take a continuous group homomorphism (or measurable factor map)
\[
\phi:\TT_{\perp}\to \TT^1=\RR/\ZZ,
\]
and let
\[
\theta_t:=\phi(u_t)\in\TT^1.
\]
Then $\theta_t$ satisfies one-dimensional rotation
\[
\theta_{t+1}=\theta_t+\alpha\quad (\mathrm{mod}\ 1),
\]
where $\alpha:=\phi(\beta)\in \TT^1$.
\end{definition}

\subsection{Periodic Approximation as Finite Resource Implementability Mechanism (Alternative Formulation)}

\begin{assumption}[Finite resource approximation principle]\label{assump:finite-resource-cf}
The observation protocol cannot realize irrational rotation parameters $\alpha$ infinitely precisely under finite computation/finite memory conditions; its implementable parameter set is produced by convergent subsequences from continued fraction truncations $\alpha_k=p_k/q_k$, where
\[
\alpha_k\to \alpha,\qquad \gcd(p_k,q_k)=1.
\]
Thus, on finite timescales the observed orbit can be approximated by closed orbits of period $q_k$; as resources increase (larger $k$) approximation precision improves, and periodic approximations form a hierarchical sequence.
\end{assumption}

\begin{remark}
The above formulation replaces ``reachable sampling'' from intuitive language with ``parameter approximation domain under resource constraint''; its mathematical content is continued fraction theory and periodic approximation of torus rotations.
\end{remark}
